\begin{table}[htbp]
  \centering
  \caption{Real hourly labor income by firm size, pooled years}
  \label{tab:descriptive-income-size}
  \small
  \begin{tabular}{lrrrrrr}
    \toprule
    Firm size & Mean & P25 & Median & P75 & S.D. & Raw premium \\
    \midrule
    Solo & 5,690 & 2,366 & 3,943 & 6,169 & 11,205 & 0.0\% \\
    2-3 & 6,464 & 3,067 & 4,784 & 6,779 & 12,543 & 13.6\% \\
    4-5 & 7,444 & 3,943 & 5,430 & 7,163 & 21,476 & 30.8\% \\
    6-10 & 8,403 & 4,647 & 5,922 & 7,854 & 16,634 & 47.7\% \\
    11-19 & 9,746 & 5,204 & 6,383 & 8,774 & 17,881 & 71.3\% \\
    20-30 & 9,972 & 5,468 & 6,669 & 9,341 & 21,069 & 75.3\% \\
    31-50 & 10,744 & 5,714 & 6,939 & 10,321 & 16,435 & 88.8\% \\
    51-100 & 11,612 & 5,809 & 7,176 & 11,255 & 18,818 & 104.1\% \\
    101+ & 14,650 & 6,166 & 8,624 & 16,169 & 20,771 & 157.5\% \\
    \bottomrule
  \end{tabular}
  \vspace{0.3em}
  \begin{minipage}{0.95\textwidth}
  \footnotesize
  Notes: Incomes are pooled over 2008--2019 and 2021--2025 and expressed as real hourly labor income in constant 2025 pesos. All statistics are weighted by GEIH expansion weights. The raw premium is the percent difference in the weighted mean relative to solo workers.
  \end{minipage}
\end{table}
